\section{Related work}
\label{sec:related}

\subsection{Vision--language pretraining and its biomedical adaptations}

Conventional recognition models are trained against a fixed, closed set of
labels chosen in advance, so every additional category requires fresh labeled
examples and usually a new output layer. Vision--language pretraining replaces
that arrangement with a joint embedding space for images and free text, learned
by contrastively aligning large corpora of image--caption pairs, so that a
category is specified at inference by writing it in natural language rather than
by having been enumerated at training time. The supervision is not free, since
captions are themselves human-written text paired with images, but it is
open-vocabulary, and that property is what makes the paradigm attractive for a
problem in which the structures of interest are not fixed in advance.

CLIP established that predicting which caption belongs to which image, over 400
million web-collected pairs, yields representations strong enough for zero-shot
transfer across more than 30 vision benchmarks \cite{radford2021clip}. SigLIP
later replaced the softmax-normalized contrastive objective with a pairwise
sigmoid loss that needs no global view of the similarity matrix, which decouples
the loss from batch size and improves results at both small and very large
batches \cite{zhai2023siglip}.

Much of the subsequent progress has been attributed to data rather than
architecture. MetaCLIP argues explicitly that CLIP's principal ingredient is its
curation procedure, and recovers CLIP-level performance by balancing a raw pool
against metadata derived from CLIP's own concept list \cite{xu2024metaclip}.
DataComp makes the same point by construction, fixing the training code and
treating the dataset as the object of study across a 12.8-billion-pair candidate
pool \cite{gadre2023datacomp}. Both results motivate the biomedical variants,
where the binding constraint is domain-matched data.

Two such adaptations are directly relevant. PMC-CLIP assembles PMC-OA, 1.6M
image--caption pairs from PubMed Central's Open Access subset, most aligned at
subfigure and subcaption granularity \cite{lin2023pmcclip}. BiomedCLIP scales
this further with PMC-15M, 15 million figure--caption pairs drawn from 4.4
million articles, and adapts the text side to the domain: a PubMedBERT encoder
with a 30k domain-specific vocabulary in place of a general-domain tokenizer,
and a context length raised to 256 tokens to accommodate the length of PMC
captions \cite{zhang2025biomedclip}. That longer context is what makes
multi-sentence anatomical prompts expressible at all, and this work builds on
BiomedCLIP for that reason.

What none of these models supplies, however, is spatial grounding, and the
reason is worth setting out as a diagnosis rather than a complaint, because
these models are behaving exactly as their objective specifies. That objective
constrains only the pooled image representation against the pooled caption
representation: the pooled vector is what is projected into the shared embedding
space and scored, and individual patch tokens enter the loss only indirectly,
through whatever the pooling operation draws from them. No gradient therefore
ever asks a particular patch token to align with the text on its own, which
means a model can minimize the objective perfectly while its patch tokens hold
representations that are useful for assembling the pooled summary yet not
individually meaningful when compared against a text embedding. Two further
properties of the architecture compound this. The projection head is fit to the
distribution of pooled representations, so pushing patch tokens through that
same head asks it to do something it was never trained for; and self-attention
in the deeper blocks mixes information globally, so a late-layer patch token
reflects image-wide context rather than the evidence at its own location, which
is precisely the locality a similarity map would need.

That the pooling stage is implicated is not a new observation. MaskCLIP
recovers a substantial amount of dense signal by reading CLIP's value projection
instead of the attention-pooled output \cite{zhou2022maskclip}, later work in
the same line pushes that recovery further \cite{yang2025clipdense}, and
methods that re-pose the objective at pixel level \cite{rao2022denseclip} or
attach a decoder trained for the purpose \cite{lueddecke2022clipseg} restore
dense behavior by other means. Read naively, then, per-patch similarity from a
CLS-trained model has no reason to carry a localization signal, and the
untrained pipeline recovers almost none: gradient saliency on the pretrained
model reaches 0.003 mIoU before any fine-tuning (\S\ref{sec:results}). The
question this work takes up is what supplies the gradient pressure the
pretraining objective never had reason to apply.

\subsection{Dense and local alignment from vision--language models}

Those repairs are worth separating out, because global alignment has been
identified as the obstacle before, and the closest work attacks it in three
distinguishable ways.

The first extracts dense correspondence from an already-trained model without
changing its weights. MaskCLIP shows that CLIP's value projection retains
substantially more spatial information than the pooled embedding suggests, and
recovers usable dense labels by reading it directly \cite{zhou2022maskclip}.
Subsequent work pushes this further for weakly supervised semantic segmentation
\cite{yang2025clipdense}, and a related line examines why CLIP's raw attention
localizes poorly and how its explainability can be repaired
\cite{li2024clipexplain}. Self-supervised vision transformers exhibit a similar
property, with DINO producing segmentation-like attention maps without any
textual supervision \cite{caron2021dino}. OrganPrompt takes the same view of
where the signal lives but trains the property in rather than recovering it,
because the downstream index requires organ maps that vanish where the organ is
absent, which is a calibration property rather than a localization one.

The second, instead of reading the signal out of frozen weights, adapts the
model to dense prediction with a task head. DenseCLIP converts image-text
matching into pixel-text matching and fine-tunes with a dense decoder
\cite{rao2022denseclip}, and CLIPSeg attaches a lightweight transformer decoder
that emits a segmentation from a text or image prompt
\cite{lueddecke2022clipseg}. Both produce better boundaries than OrganPrompt,
but both also place a trained decoder between the prompt and the mask, which is
precisely the attribution problem described below.

The third builds the local property in during pretraining rather than adding it
afterwards, and GLoRIA is the closest medical antecedent: it learns global and
local representations jointly by matching image sub-regions to individual words
in radiology reports \cite{huang2021gloria}. The objective here is the same in
spirit, but it differs in supervision and in target, because GLoRIA attends over
report words to build localized representations for label-efficient recognition,
whereas we supervise each of the 196 patch tokens directly against a downsampled
organ mask, so that the resulting per-patch similarity is itself the
segmentation rather than an intermediate representation.

\subsection{Text-driven medical image segmentation}

Within text-driven medical segmentation itself, the closest methodological
relative is MedCLIP-SAMv2, which shares both the BiomedCLIP backbone and the
premise that a text prompt can drive segmentation
\cite{koleilat2025medclipsamv2}, extending an earlier version of the same
framework \cite{koleilat2024medclipsam}. It differs in where the mask comes from:
saliency derived from the vision--language model is used to prompt SAM
\cite{kirillov2023sam}, which produces the final boundary, in a lineage that
includes MedSAM's adaptation of SAM to medical images \cite{ma2024medsam}.

Delegating mask formation to a second model, as that lineage does, is the
stronger engineering choice and generally yields better boundaries, but it also
changes what the pipeline can be asked. When a refiner sits between the prompt
and the output, a failure cannot be attributed cleanly to the text, because the
refiner may repair a bad localization or corrupt a good one. Reading patch
similarity directly, as here, gives up boundary quality and gains attribution in
exchange: the caption ablation in \S\ref{sec:results:ablation} is only
interpretable because nothing stands between the prompt and the mask. We regard
that trade as the substantive difference between the two approaches, not the
scores in Table~\ref{tab:weaksup}, which are computed on different datasets and
organ sets.

\subsection{Weak supervision, stratified by annotation cost}

That lineage and OrganPrompt are both weakly supervised in the broad sense,
which is exactly the problem: ``weakly supervised'' covers a wide range of
supervision strength, and comparing across that range without saying so is
misleading.
Point-supervised methods receive one clicked pixel inside each structure;
scribble-supervised methods receive a stroke through it. Recent work in this
tier reports substantial gains
over prior point- and scribble-supervised baselines on BTCV and CHAOS
\cite{jiang2024probabilistic}. Multi-organ abdominal CT benchmarks such as AMOS
\cite{ji2022amos} and WORD \cite{luo2022word} are the usual evaluation targets in
this tier, which is one reason no row of Table~\ref{tab:weaksup} shares a dataset
with the evaluation here.

What point and scribble annotation have in common is that both answer ``where''
for each training image. Image-level text supervision does not answer it at all,
which is the entire problem being solved here, so these methods are context
rather than baselines. The strictly comparable tier is instead image-level
supervision, where weakly supervised CAM-based lesion segmentation in PET-CT
reports a median Dice of 0.47 against 0.72 for a fully supervised U-Net
\cite{ronneberger2015unet} on the same task \cite{fruh2021petct}, the fully
supervised tier in which transformer architectures such as TransUNet
\cite{chen2021transunet} are now a common alternative. That ratio, roughly two
thirds of supervised performance recovered from image-level labels alone, is the
expectation this work should be read against.

\subsection{Class activation maps and post-hoc localization}

The image-level tier has an established localization technique of its own, and
it starts from a classifier rather than from a vision--language model: the CAM
family. CAM recovers localization from a classifier's final-layer activations
\cite{zhou2016cam}, and Grad-CAM generalizes this to arbitrary architectures by
weighting activations with backpropagated gradients \cite{selvaraju2017gradcam},
with later refinements such as Grad-CAM++ improving multi-instance localization
\cite{chattopadhay2018gradcampp}. What every member of the family has in common
is that it explains where a trained classifier looked.

That shared framing is what makes the family a poor fit here, for a reason worth
stating plainly rather than asserting. A CAM explains a decision the network
already makes, and a contrastively pretrained vision--language model makes no
per-organ decision to explain, because it produces an embedding whose alignment
with a text embedding is a similarity rather than a class logit. Applying
gradient-based localization to it consequently recovers almost nothing, which is
what the 0.003 starting point in Table~\ref{tab:progression} measures. Because
there is no decision to explain, this work trains the spatial property directly
instead of attempting to extract it post hoc, and the name reflects that: the
prompt does the localizing, and no activation map is involved.

\subsection{Positioning}

Those observations locate this work. It sits in the image-level tier and differs
from its nearest neighbors in three respects, the first of which the preceding
subsection has already anticipated. It trains patch-level text--image alignment
as an explicit objective instead of extracting localization post hoc from a
classifier. It reads the resulting similarity map directly, with no second-stage
refiner, which preserves the ability to attribute failures to the prompt. And it
carries the same frozen model into a downstream
clinical question, the bone-metastasis proximity index of
\S\ref{sec:results:index}, on a dataset and acquisition protocol the model never
saw during training.

One point of provenance completes that picture. The underlying idea was
validated on the OrganAMNIST subset of MedMNIST v2 \cite{yang2023medmnist}
before the move to full-resolution CT; the results here are on the latter only.
